\section{Modules of finite presentation}
\label{section-finite-presentation}

\noindent
Here is the definition.

\begin{definition}
\label{definition-finite-presentation}
Let $(X, \mathcal{O}_X)$ be a ringed space.
Let $\mathcal{F}$ be a sheaf of $\mathcal{O}_X$-modules.
We say that $\mathcal{F}$ is of {\it finite presentation}
if for every point $x \in X$ there exists an open neighbourhood
$x\in U \subset X$, and  $n, m \in \mathbf{N}$ such that $\mathcal{F}|_U$
is isomorphic to the cokernel of a map
$$
\bigoplus\nolimits_{j = 1, \ldots, m}
\mathcal{O}_U
\longrightarrow
\bigoplus\nolimits_{i = 1, \ldots, n}
\mathcal{O}_U
$$
\end{definition}

\noindent
This means that $X$ is covered by open sets $U$
such that $\mathcal{F}|_U$ has a {\it presentation}
of the form
$$
\bigoplus\nolimits_{j = 1, \ldots, m}
\mathcal{O}_U
\longrightarrow
\bigoplus\nolimits_{i = 1, \ldots, n}
\mathcal{O}_U
\to
\mathcal{F}|_U
\to
0.
$$
Here presentation signifies that the displayed
sequence is exact. In other words
\begin{enumerate}
\item for every point $x$ of $X$ there exists
an open neighbourhood such that $\mathcal{F}|_U$
is generated by finitely many global sections, and
\item for a suitable choice of these sections
the kernel of the associated surjection is also
generated by finitely many global sections.
\end{enumerate}

\begin{lemma}
\label{lemma-finite-presentation-quasi-coherent}
Let $(X, \mathcal{O}_X)$ be a ringed space.
Any $\mathcal{O}_X$-module of finite presentation
is quasi-coherent.
\end{lemma}

\begin{proof}
Immediate from definitions.
\end{proof}

\begin{lemma}
\label{lemma-cokernel-finite-finite-presentation}
Let $(X,\mathcal{O}_X)$ be a ringed space. Let $\mathcal{F}$ be a
finitely presented $\mathcal{O}_X$-module. Let
$\varphi : \mathcal{G} \to \mathcal{F}$ be a morphism of
$\mathcal{O}_X$-modules. If $\mathcal{G}$ is finite type, then
$\Coker(\varphi)$ is finitely presented.
\end{lemma}

\begin{proof}
Locally on $X$, we can write
$\mathcal{F} = \mathcal{O}_X^{\oplus n} / \mathcal{M}$, where
$\mathcal{M} \subset \mathcal{O}_X^{\oplus n}$ is a finite type
$\mathcal{O}_X$-submodule. Therefore
$\Im(\varphi) = \mathcal{N} / \mathcal{M}$, where
$\mathcal{N} \subset \mathcal{O}_X^{\oplus n}$ is a
$\mathcal{O}_X$-submodule containing $\mathcal{M}$.
The $\mathcal{O}_X$-module $\Im(\varphi)$ is of finite type
because $\mathcal{G} \to \Im(\varphi)$ is surjective and
$\mathcal{G}$ is finite type. By Lemma \ref{lemma-extension-finite-type},
$\mathcal{N}$ is finite type. Thus
$\Coker(\varphi) = \mathcal{O}_X^{\oplus n} / \mathcal{N}$
is finitely presented.
\end{proof}

\begin{lemma}
\label{lemma-kernel-surjection-finite-free-onto-finite-presentation}
Let $(X, \mathcal{O}_X)$ be a ringed space.
Let $\mathcal{F}$ be an $\mathcal{O}_X$-module of finite presentation.
\begin{enumerate}
\item If $\psi : \mathcal{O}_X^{\oplus r} \to \mathcal{F}$ is a surjection,
then $\Ker(\psi)$ is of finite type.
\item If $\theta : \mathcal{G} \to \mathcal{F}$ is surjective with
$\mathcal{G}$ of finite type, then $\Ker(\theta)$ is of finite type.
\end{enumerate}
\end{lemma}

\begin{proof}
Proof of (1).
Let $x \in X$. Choose an open neighbourhood $U \subset X$ of $x$
such that there exists a presentation
$$
\mathcal{O}_U^{\oplus m}
\xrightarrow{\chi}
\mathcal{O}_U^{\oplus n}
\xrightarrow{\varphi}
\mathcal{F}|_U
\to
0.
$$
Let $e_k$ be the section generating the $k$th factor of
$\mathcal{O}_X^{\oplus r}$. For every $k = 1, \ldots, r$
we can, after shrinking $U$ to a small neighbourhood of $x$,
lift $\psi(e_k)$ to a section $\tilde e_k$ of
$\mathcal{O}_U^{\oplus n}$ over $U$. This gives a morphism
of sheaves $\alpha : \mathcal{O}_U^{\oplus r} \to \mathcal{O}_U^{\oplus n}$
such that $\varphi \circ \alpha = \psi$.
Similarly, after shrinking $U$, we can find a morphism
$\beta : \mathcal{O}_U^{\oplus n} \to \mathcal{O}_U^{\oplus r}$
such that $\psi \circ \beta = \varphi$. Then the map
$$
\mathcal{O}_U^{\oplus m} \oplus
\mathcal{O}_U^{\oplus r}
\xrightarrow{\beta \circ \chi, 1 - \beta \circ \alpha}
\mathcal{O}_U^{\oplus r}
$$
is a surjection onto the kernel of $\psi$.

\medskip\noindent
To prove (2) we may locally choose a surjection
$\eta : \mathcal{O}_X^{\oplus r} \to \mathcal{G}$.
By part (1) we see $\Ker(\theta \circ \eta)$ is of finite type.
Since $\Ker(\theta) = \eta(\Ker(\theta \circ \eta))$ we win.
\end{proof}

\begin{lemma}
\label{lemma-pullback-finite-presentation}
Let $f : (X, \mathcal{O}_X) \to (Y, \mathcal{O}_Y)$
be a morphism of ringed spaces.
The pullback $f^*\mathcal{G}$ of a module of finite presentation
is of finite presentation.
\end{lemma}

\begin{proof}
Exactly the same as the proof of Lemma \ref{lemma-pullback-quasi-coherent}
but with finite index sets.
\end{proof}

\begin{lemma}
\label{lemma-quasi-coherent-limit-finite-presentation}
Let $(X, \mathcal{O}_X)$ be a ringed space.
Set $R = \Gamma(X, \mathcal{O}_X)$.
Let $M$ be an $R$-module.
The $\mathcal{O}_X$-module $\mathcal{F}_M$ associated to $M$
is a directed colimit of finitely presented $\mathcal{O}_X$-modules.
\end{lemma}

\begin{proof}
This follows immediately from
Lemma \ref{lemma-construct-quasi-coherent-sheaves}
and the fact that any module is a directed colimit
of finitely presented modules, see
Algebra, Lemma \ref{algebra-lemma-module-colimit-fp}.
\end{proof}

\begin{lemma}
\label{lemma-finite-presentation-stalk-free}
Let $(X, \mathcal{O}_X)$ be a ringed space. Let $\mathcal{F}$ be
a finitely presented $\mathcal{O}_X$-module. Let $x \in X$ such that
$\mathcal{F}_x \cong \mathcal{O}_{X, x}^{\oplus r}$. Then there exists
an open neighbourhood $U$ of $x$ such that
$\mathcal{F}|_U \cong \mathcal{O}_U^{\oplus r}$.
\end{lemma}

\begin{proof}
Choose $s_1, \ldots, s_r \in \mathcal{F}_x$ mapping to a basis of
$\mathcal{O}_{X, x}^{\oplus r}$ by the isomorphism. Choose an open
neighbourhood $U$ of $x$ such that $s_i$ lifts to $s_i \in \mathcal{F}(U)$.
After shrinking $U$ we see that the induced map
$\psi : \mathcal{O}_U^{\oplus r} \to \mathcal{F}|_U$ is surjective
(Lemma \ref{lemma-finite-type-surjective-on-stalk}).
By Lemma \ref{lemma-kernel-surjection-finite-free-onto-finite-presentation}
we see that $\Ker(\psi)$ is of finite type.
Then $\Ker(\psi)_x = 0$ implies that $\Ker(\psi)$ becomes zero
after shrinking $U$ once more (Lemma \ref{lemma-finite-type-stalk-zero}).
\end{proof}





